% -----------------------------------------------------------
\section{Buffet}
% -----------------------------------------------------------
	\label{BUFFET}
	
% % ----------------------
% \subsection{See also}
% % ----------------------

% % ----------------------
% \subsection{1828 American Dictionary of the English Language} % *1828
% % ----------------------
	% \label{BUFFET-1828}

	% \begin{compactitem}
		% \item
	% \end{compactitem}

% % ----------------------
% \subsection{Oxford English Dictionary} % *OED
% % ----------------------
	% \label{BUFFET-OED}

	% \begin{compactitem}
		% \item[]
	% \end{compactitem}
	
% % ----------------------
% \subsection{Restoration Lexicon} % *RL *RESTORATION LEXICON
% % ----------------------
	% \label{BUFFET-OED}

	% \begin{compactitem}
		% \item[]
	% \end{compactitem}

% ----------------------
\subsection{Scriptural Usage} % *SCRIPTURES
% ----------------------
	\label{BUFFET-SCRIPTURES}

	% % -----
	% \subsubsection{Old Testament} % *OT
	% % -----
		% \label{BUFFET-OT}
	
		% --
		\paragraph{KJV}
		% --
			\label{BUFFET-OT-KJV}
	
			% \begin{tabular}{ll}
				% Total occurrences in the OT & 
			% \end{tabular}
			
			\begin{compactdesc}
				\item[{\hyperref[MATTHEW-5-25]{Matthew 5:25--26}}]
					\begin{compactitem}
						\item These verses don't use the term ``buffet,'' but it may be a way of understanding a similar idea.
					\end{compactitem}
			\end{compactdesc}
			
		% % --
		% \paragraph{Hebrew Bible}
		% % --
			% \label{BUFFET-HEBREW}
		
			% %
			% \subparagraph{\texorpdfstring{\he{sample word}}{}} % paste actual hebrew text here
			% %
				% \label{PAGA} % whatever is in step bible without periods
			
				% \begin{compactdesc}
					% \item[HALOT , PDF ] \textbf{} % Use the 1994 PDF
						% \begin{compactitem}
							% \item[1]
						% \end{compactitem}
					% \item[BDB , PDF ] \textbf{}
						% \begin{compactitem}
							% \item[1]
						% \end{compactitem}
					% \item[DCH PDF ] \textbf{} % don't include the actual page number, they repeat from volume to volume
						% \begin{compactitem}
							% \item[1]
						% \end{compactitem}
				% \end{compactdesc}
				
				% \begin{compactdesc}
					% \item[Some OT uses] \textbf{}
						% \begin{compactdesc}
							% \item[]
						% \end{compactdesc}
				% \end{compactdesc}
			
		% % --
		% \paragraph{Septuagint}
		% % --
			% \label{BUFFET-LXX}
		
			% %
			% \subparagraph{\texorpdfstring{\gr{sample word}}{}}
			% %
		
	% -----
	\subsubsection{New Testament} % *NT
	% -----
		\label{BUFFET-NT}
	
		% --
		\paragraph{KJV}
		% --
			\label{BUFFET-NT-KJV}		
	
			% \begin{tabular}{ll}
				% Total occurrences in the NT & 
			% \end{tabular}
			
			\begin{compactdesc}
				\item[{\hyperref[2CORINTHIANS-12-7]{2 Corinthians 12:7}}]
			\end{compactdesc}
			
		% --
		\paragraph{Greek New Testament} % *GREEK
		% --
			\label{BUFFET-GREEK}
		
			%
			\subparagraph{\texorpdfstring{\gr{κολαφίζω}}{}}
			%
				\label{KOLAPHIZO}
			
				\begin{compactdesc}
					\item[BDAG 492a] \textbf{}
						\begin{compactitem}
							\item[1] \textbf{to strike sharply, esp. with the hand, \textit{strike with the fist, beat, cuff}}
							\item[2] \textbf{to cause physical impairment, \textit{torment}}, fig. extension of 1, of painful attacks of an illness, described as a physical beating by a messenger of Satan 2 Cor 12:7. The data for a scientific diagnosis are few, and it is not surprising that a variety of views, characterized by much guesswork, have been held: epilepsy, hysteria, periodic depression, headaches/severe eye trouble, malaria, leprosy, an impediment in speech
								\begin{compactitem}
									\item It does mention very briefly that many people formely favored an interpretation of this term in 2 Corinthians 12 where it had more to do with inward temptations.
								\end{compactitem}
						\end{compactitem}
					% \item[CGL , PDF ] \textbf{}
						% \begin{compactitem}
							% \item 
						% \end{compactitem}
					% \item[LSJ , PDF ] \textbf{}
						% \begin{compactitem}
							% \item 
						% \end{compactitem}
				\end{compactdesc}
				
				% \begin{tabular}{ll}
					% Total occurrences in the NT & 
				% \end{tabular}
				
				% \begin{compactdesc}
					% \item[Some NT uses] \textbf{}
						% \begin{compactdesc}
							% \item[]
						% \end{compactdesc}
				% \end{compactdesc}
				
			% %
			% \subparagraph{TDNT: \texorpdfstring{\gr{sample word}}{} (3:536--556)}
			% %
				% \label{KATALLAGE-TDNT}
		
	% % -----
	% \subsubsection{Book of Mormon} % *BOM
	% % -----
		% \label{BUFFET-BOM}
	
		% \begin{tabular}{ll}
			% Total occurrences in the BOM & 
		% \end{tabular}
		
		% \begin{compactdesc}
			% \item[]
		% \end{compactdesc}
		
	% -----
	\subsubsection{Doctrine and Covenants} % *DC
	% -----
		\label{BUFFET-DC}
	
		% \begin{tabular}{ll}
			% Total occurrences in the DC & 
		% \end{tabular}
		
		\begin{compactdesc}
			\item[{\hyperref[DC78-12]{DC 78:12}}]
			\item[{\hyperref[DC82-21]{DC 82:21}}]
			\item[{\hyperref[DC132-26]{DC 132:26}}]
		\end{compactdesc}
		
	% % -----
	% \subsubsection{Pearl of Great Price} % *PGP
	% % -----
		% \label{BUFFET-PGP}
	
		% \begin{tabular}{ll}
			% Total occurrences in the PGP & 
		% \end{tabular}
		
		% \begin{compactdesc}
			% \item[]
		% \end{compactdesc}
		
% ----------------------
\subsection{Key Phrases} % *KEYPHRASES *PHRASES
% ----------------------
	\label{BUFFET-PHRASES}

	\begin{compactdesc}
		\item[buffetings of Satan] \textbf{6} | \hyperref[2CORINTHIANS-12-7]{2 Corinthians 12:7}; \hyperref[DC78-12]{DC 78:12}; \hyperref[DC82-21]{DC 82:21}; \hyperref[DC104-9]{DC 104:9}; \hyperref[DC104-10]{DC 104:10}; \hyperref[DC132-26]{DC 132:26}
			\begin{compactdesc}
				\item[Bible anchor verses] \phantom{.}
					\begin{compactdesc}
						\item[{\hyperref[2CORINTHIANS-12-7]{2 Corinthians 12:7}}] \gr{καὶ τῇ ὑπερβολῇ τῶν ἀποκαλύψεων. διὸ ἵνα μὴ ὑπεραίρωμαι, ἐδόθη μοι σκόλοψ τῇ σαρκί, ἄγγελος Σατανᾶ, ἵνα με κολαφίζῃ, ἵνα μὴ ὑπεραίρωμαι.}
					\end{compactdesc}
				\item[Commentary] ---
			\end{compactdesc}
	\end{compactdesc}
	
% % ----------------------
% \subsection{Bible Dictionaries} % *DICTIONARIES
% % ----------------------
	% \label{BUFFET-BD}

% % ----------------------
% \subsection{Restoration Usage} % *RESTORATION
% % ----------------------
	% \label{BUFFET-RESTORATION}

	% % -----
	% \subsubsection{Joseph Smith} % *JS
	% % -----
		% \label{BUFFET-JS}
	
		% \begin{compactdesc}
			% \item[{\hyperref[KEY]{Date/Source}}]
		% \end{compactdesc}
		
	% % -----
	% \subsubsection{Temple Ordinances}
	% % -----
		% \label{BUFFET-TEMPLE}
	
		% \begin{compactdesc}
			% \item[{\hyperref[KEY]{Date/Source}}]
		% \end{compactdesc}
	
	% % -----
	% \subsubsection{Brigham Young} % *BY
	% % -----
		% \label{BUFFET-BY}
	
		% \begin{compactdesc}
			% \item[{\hyperref[KEY]{Date/Source}}]
		% \end{compactdesc}
	
% ----------------------
\subsection{Commentary and Studies} % *COMMENTARY *STUDIES
% ----------------------
	\label{BUFFET-COMMENTARY}

	% % -----
	% \subsubsection{Key Points}
	% % -----
		% \label{BUFFET-KEYS}
	
		% \begin{compactitem}
			% \item
		% \end{compactitem}
		
	% -----
	\subsubsection{Ken Openshaw's discussion of the ``buffetings of Satan'' (January 2023)}
	% -----
		\label{BUFFET-openshaw}
		
		Our objective in coming to this Earth is to become like God in every way: characteristics, eternal body, glory, power, priesthood, fatherhood, and etcetera. We call becoming like God in every way, Eternal Life. For many years, I considered the Buffetings of Satan only in a negative way. I viewed it as a punishment inflicted upon those who disobeyed. I now view the Buffetings of Satan as more of a gift, a blessing provided to those who are in need of a certain experience to overcome weaknesses or be exposed to something that they have not been exposed to but is required in order to become like God in every way. 

		Christ's work and his glory is to bring to pass the Immortality and Eternal Life of man. (Moses 1:39) He accomplishes His part through (1) His Resurrection; (2) His Atonement and (3) His Example. The resurrection provides every individual who comes to this earth with an eternal body regardless of anything we do. The atonement gifts to those of us who have faith in Jesus Christ and accept His atonement the ability to be forgiven of any act we do (through repentance) or is done to us (freely) that separates us from God. His example gives us the ``road map'' of what we must do to become like Jesus Christ and therefore like the Father. This example shows us what Christ did to become like the Father and therefore what we must do to become like the Father.

		Christ could not have become like the Father if He did not do everything that His Father did. (John 5:19) This means that Christ had to control his earthly body, confront Satan, love every son and daughter born to the Father and perform the atonement, among other things. It is therefore logical that we cannot become like the Father if we do not do everything that Christ did.

		We have two barriers in this life that limit our doing everything that Christ did: (1) we never have the opportunity and (2) we are incapable of doing what Christ did through physical or mental limitations or spiritual weaknesses. God's plan permits us to overcome these barriers through the Atonement and the Buffetings of Satan. 

		The Atonement cleanses us from any act that separates us from God. Regardless of how many times we sincerely ask for cleansing, even for repeated offenses, we receive it. (Moroni 6:8) By being cleansed, we are empowered to try again to change our behavior or overcome that situation that we cannot control. The Atonement does not, however, take away that weakness. We are the ones who have to work at resolving that issue that separates us from God. The scriptures call this process of becoming perfected ``...paying the penalty of our transgressions.'' (DC 138:59) John described the Savior's efforts in overcoming the weaknesses of man as Christ's obtaining His fullness by receiving Grace for Grace and then continuing in Grace to Grace. (DC 93:11--12) Christ received grace (forgiveness) every time he did something that separated Him from God and since He never repeated those acts, he continued from grace to the next time He needed grace.

		The Buffetings of Satan is that process that permits Satan to tempt, challenge, confront and test us with essential experiences that are required to become like the Father. Since we all have different life experiences, The Buffetings of Satan is one of the blessings that the Father provides as part of His Plan that permits us to have those essential experiences that will permit us to become like Him. 

		The Buffetings of Satan can be experienced while we are on this earth or in the Spirit World. (DC 132:26) Therefore every experience that we lack to become a God will be provided as long as we desire to continue progressing.

		Because Satan's ability to test us can overpower us if we are not ready for the challenge, Satan is generally limited on what he is permitted to do to us. These limitations are controlled by God and as we request and are ready, He eases these restrictions so we can progress through our personal development. Eventually, if we continue our progression, we will be given the opportunity to confront Satan in his full power and capacity as well as confront every weakness we need to overcome.